\section{Compromessi e Sviluppi
Futuri}\label{compromessi-e-sviluppi-futuri}

Nonostante i vantaggi prestazionali e l'isolamento garantito
dall'impiego congiunto di ZFS, Incus e OVN, l'architettura di BlueAgent
presenta alcuni compromessi progettuali e delinea diverse opportunità di
estensione futura.

\subsection{Compromessi
Architetturali}\label{compromessi-architetturali}

\begin{enumerate}
\def\labelenumi{\arabic{enumi}.}
\item
  \textbf{MTU Blackhole per il traffico UDP:}\\
  Il traffico tra i partecipanti e i container transita attraverso il
  tunnel WireGuard, caratterizzato da una \(\text{MTU}\) di
  \(1420\text{ byte}\). Sebbene il problema della frammentazione sia
  risolto per il traffico TCP tramite il clamping dell'MSS in
  \passthrough{\lstinline!nftables!}, i pacchetti UDP non dispongono di
  un meccanismo di negoziazione analogo. Di conseguenza, datagrammi UDP
  generati con flag \emph{Don't Fragment} (DF) e dimensione standard
  Ethernet (\(1500\text{ byte}\)) vengono scartati dal kernel senza
  recapito.
\item
  \textbf{Limite di scalabilità dello schema di indirizzamento:}\\
  La mappatura deterministica \passthrough{\lstinline!10.k.x.x!} associa
  l'identificatore del team \(k\) al secondo ottetto dell'indirizzo IPv4
  (\(1 \le k \le 254\)). Questo vincolo delimita la capacità del sistema
  a un massimo di 254 team concorrenti per singola istanza, richiedendo
  una segmentazione a subnet più estese (o l'adozione di IPv6) per
  scenari di scala superiore.
\item
  \textbf{Confinamento del locking al singolo host:}\\
  Il sistema di sincronizzazione gerarchico
  \passthrough{\lstinline!TwoTierLock!} impiega
  \passthrough{\lstinline!asyncio.Lock!} a livello di processo e
  \passthrough{\lstinline!fcntl.flock!} a livello di filesystem locale.
  Questa soluzione garantisce consistenza su macchine multi-core
  mononodo, ma non è estendibile a cluster distribuiti. Scalare
  l'orchestrazione su più server fisici richiederebbe l'introduzione di
  un gestore di lock distribuito (es. Redis o etcd).
\item
  \textbf{Trade-off tra flessibilità e dipendenze esterne:}\\
  Per mantenere l'infrastruttura minimale e immediata nel deployment, si
  è scelto di non introdurre middleware complessi quali Kubernetes o
  Ansible. Tale approccio elimina complessità di configurazione, ma
  vincola la logica di rete delle challenge alle funzionalità esposte
  internamente da BlueAgent, rendendo meno agevole l'implementazione di
  tipi di challenge eterogenee.
\end{enumerate}

\subsection{Sviluppi Futuri}\label{sviluppi-futuri}

\begin{enumerate}
\def\labelenumi{\arabic{enumi}.}
\item
  \textbf{Interfaccia da riga di comando
  (\passthrough{\lstinline!bluecmd!}):}\\
  Sviluppo di una CLI dedicata per fornire ad amministratori e creatori
  di challenge un tool ufficiale ed affidabile per gestire il ciclo di
  vita delle istanze.
\item
  \textbf{Monitoraggio dello SLA e Service Health Checking:}\\
  Integrazione di verifiche periodiche di disponibilità dei servizi di
  produzione. Il controllo continuo dello SLA (\emph{Service Level
  Agreement}) è necessario in una competizione di Difesa per penalizzare
  chi cerca di barare spegnendo il servizio.
\item
  \textbf{Configurazione dichiarativa avanzata delle challenge:}\\
  Estensione del parser YAML per supportare scenari multi-tier complessi
  (es. router virtuali dedicati) e configurare per ogni container le
  interfacce di rete.
\end{enumerate}
